\begingroup
\renewcommand{\arraystretch}{1.45}
\fancyhead[L]{\footnotesize ONE HUNDRED THEOREMS ON THE NEW TRANSCENDENTAL FUNCTIONS.}

\Needspace{32\baselineskip}
\par\medskip\Needspace{5\baselineskip}\pdfbookmark[1]{Art. 4}{transscendentes-iv-art-4}\centerline{4.}\nobreak\smallskip
If $n$ is allowed to increase to infinity, the series and products of the second and third theorems become infinite series and products. In this form, the \textsc{fourth theorem} is
\[
\begin{gathered}
1+x(y+\frac1y)+x^4(yy+\frac1{yy})+x^9(y^3+\frac1{y^3})+\text{etc. in inf.}\\[8pt]
=(1+xy)(1+\frac xy)(1+x^3y)(1+\frac{x^3}y)\\[8pt]\quad{}(1+x^5y)(1+\frac{x^5}y)\mathinner{.\,.\,.}\\[8pt]
\quad\times(1-xx)(1-x^4)(1-x^6)(1-x^8)\mathinner{.\,.}
\end{gathered}
\]

\Needspace{19\baselineskip}
and \textsc{the fifth}
\[
\begin{gathered}
x^{\frac14}(y^{\frac12}+y^{-\frac12})+x^{\frac94}(y^{\frac32}+y^{-\frac32})\\[8pt]\quad{}+x^{\frac{25}4}(y^{\frac52}+y^{-\frac52})+\text{etc. in inf.}\\[8pt]
=x^{\frac14}(y^{\frac12}+y^{-\frac12}).(1+xxy)(1+x^4y)(1+x^6y)\mathinner{.\,.}\\[8pt]
\quad\times(1+\frac{xx}y)\quad(1+\frac{x^4}y)\quad(1+\frac{x^6}y)\mathinner{.\,.}\\[8pt]
\quad\times(1-xx)\quad(1-x^4)\quad(1-x^6)\mathinner{.\,.}
\end{gathered}
\]

\SourceLeaf{474}


\Needspace{28\baselineskip}
\par\medskip\Needspace{5\baselineskip}\pdfbookmark[1]{Art. 5}{transscendentes-iv-art-5}\centerline{5.}\nobreak\smallskip
The functions expanded into infinite products by the fourth and fifth theorems are of great importance, and it will be useful to denote them here by special function symbols. We therefore write
\[
P(x,y)=1+x(y+y^{-1})+x^4(yy+y^{-2})+x^9(y^3+y^{-3})+\text{etc.}
\]
\[
\begin{gathered}R(x,y)=x^{\frac14}(y^{\frac12}+y^{-\frac12})+x^{\frac94}(y^{\frac32}+y^{-\frac32})\\[8pt]\quad{}+x^{\frac{25}4}(y^{\frac52}+y^{-\frac52})+\text{etc.}\end{gathered}
\]
and at the same time
\[
Q(x,y)=1-x(y+y^{-1})+x^4(yy+y^{-2})-x^9(y^3+y^{-3})+\text{etc.}
\]

\Needspace{25\baselineskip}
so that $Q(x,y)=P(-x,y)=P(x,-y)$. The simplest value that can be assigned to $y$ is $1$, and since the corresponding values of our function are of especially great importance and will occur frequently, we write, for brevity, $P(x,1),Q(x,1),R(x,1)$ simply as $Px,Qx,Rx$. We also note that when an exponent refers only to the argument of a function, this is indicated by parentheses, as in $P(x^3)$; without parentheses, it must always be understood to refer to the function, so that $Px^3$ means the same as $(Px)^3$. Thus
\[
\begin{gathered}
Px=1+2x+2x^4+\mathord{.}\\[8pt]
Qx=1-2x+2x^4-\mathord{.}\\[8pt]
Rx=2x^{\frac14}+2x^{\frac94}+2x^{\frac{25}4}+\mathord{.}
\end{gathered}
\]

\Needspace{12\baselineskip}
Finally, in this memoir we shall use the function symbol $F$ to express the infinite product
\[
Fx=(1-x)(1-xx)(1-x^3)(1-x^4)\mathinner{.\,.\,.}
\]

\Needspace{22\baselineskip}
In this notation, the last two theorems take the form
\[
\begin{gathered}\text{4.}\qquad P(x,y)=(1+xy)(1+\frac xy)(1+x^3y)(1+\frac{x^3}y)\\[8pt]\quad{}(1+x^5y)(1+\frac{x^5}y)\mathinner{.\,.\,.} Fxx\end{gathered}
\]
\[
\begin{gathered}\text{5.}\qquad R(x,y)=x^{\frac14}(y^{\frac12}+y^{-\frac12})(1+xxy)(1+\frac{xx}y)\\[8pt]\quad{}(1+x^4y)(1+\frac{x^4}y)\mathinner{.\,.\,.} Fxx\end{gathered}
\]
And thus, obviously
\[
\begin{gathered}\text{6.}\qquad Q(x,y)=(1-xy)(1-\frac xy)(1-x^3y)(1-\frac{x^3}y)\\[8pt]\quad{}(1-x^5y)(1-\frac{x^5}y)\mathinner{.\,.\,.} Fxx\end{gathered}
\]





\SourceLeaf{475}


\Needspace{22\baselineskip}
further, setting $y=1$
\[
\begin{gathered}
7.\quad Px=(1+x)^2(1-xx)(1+x^3)^2(1-x^4)\\[8pt]\quad{}(1+x^5)^2(1-x^6)\mathinner{.\,.\,.}\\[8pt]
8.\quad Qx=(1-x)^2(1-xx)(1-x^3)^2(1-x^4)\\[8pt]\quad{}(1-x^5)^2(1-x^6)\mathinner{.\,.\,.}\\[8pt]
9.\quad Rx=2x^{\frac14}(1+xx)^2(1-xx)\\[8pt]\quad{}(1+x^4)^2(1-x^4)(1+x^6)^2(1-x^6)\mathinner{.\,.}
\end{gathered}
\]

\Needspace{18\baselineskip}
Substituting here $1+x=\dfrac{1-xx}{1-x}$, $1+x^3=\dfrac{1-x^6}{1-x^3}$, etc., transforms these expressions into the following
\[
\begin{gathered}
10.\quad Px=\frac{(Fxx)^5}{(Fx)^2(Fx^4)^2}\\[8pt]
11.\quad Qx=\frac{(Fx)^2}{Fxx}\\[8pt]
12.\quad Rx=2x^{\frac14}\frac{(Fx^4)^2}{Fxx}
\end{gathered}
\]

\Needspace{29\baselineskip}
from this there further follows
\[
\begin{gathered}
13.\quad Px.Qx=(Qxx)^2\\[8pt]
14.\quad Px.Rx=\tfrac12(R\mathord{\surd}{x})^2\qquad\text{or equivalently}
\end{gathered}
\]
\[
Pxx.Rxx=\tfrac12(Rx)^2,\qquad Rxx=\frac{(Rx)^2}{2Pxx}
\]
\[
Qxx(Rx)^2=4x^{\frac12}.(Fx^4)^3\qquad\text{thus}
\]
\[
\begin{gathered}15.\quad Fx\\[8pt]\quad{}=\mathord{{}^{3}\!\surd}{\frac{Q(x^{\frac12})(Rx^{\frac14})^2}{4x^{\frac18}}}
\\[8pt]\quad{}=\mathord{{}^{3}\!\surd}{\frac{P(x^{\frac12})Q(x^{\frac12})R(x^{\frac12})}{2x^{\frac18}}}
\\[8pt]\quad{}=\mathord{{}^{3}\!\surd}{\frac{(Qx)^2R(x^{\frac12})}{2x^{\frac18}}}\end{gathered}
\]

\Needspace{16\baselineskip}
further
\[
\begin{gathered}
16.\quad Px+Qx=2P(x^4)\\[8pt]
17.\quad Px-Qx=2R(x^4)
\end{gathered}
\]
Thus, by multiplication, using 14
\[
18.\quad (Px)^2-(Qx)^2=2(Rxx)^2
\]

\Needspace{14\baselineskip}
If, further, $i$ denotes the imaginary quantity $\mathord{\surd}{-1}$, then
\[
\begin{gathered}
19.\quad Px+iQx=(1+i)Q(ix)\\[8pt]
20.\quad Px-iQx=(1-i)P(ix)
\end{gathered}
\]

\Needspace{12\baselineskip}
Thus, by multiplication

\SourceLeaf{476}

\[
21.\quad (Px)^2+(Qx)^2=2(Pxx)^2
\]
and from multiplying 18 and 21, with the use of 14
\[
22.\quad (Px)^4-(Qx)^4=(Rx)^4
\]

\Needspace{29\baselineskip}
It is therefore seen that $(Pxx)^2$ is the arithmetic mean between $(Px)^2$ and $(Qx)^2$, and since, by 13, the quantity $(Qxx)^2$ represents the geometric mean between these same quantities, and since (\textit{Theor. attract. el. p.}), if two series of quantities
\[
\begin{array}{llll}
m,&m',&m'',&m'''\mathinner{.\,.\,.}\\[8pt]
n,&n',&n'',&n'''\mathinner{.\,.\,.}
\end{array}
\]
are connected so that $m^{(\lambda)}$ is always the arithmetic and $n^{(\lambda)}$ the geometric mean between $m^{(\lambda-1)}$ and $n^{(\lambda-1)}$, their common limit is called the arithmetic–geometric mean of $m$, $n$, or of any corresponding pair of quantities in the two series, the following \textsc{Theorem} (23), of the greatest importance, results:

\textit{The arithmetic–geometric mean between $(Px)^2$ and $(Qx)^2$ is always $=1$.}


\Needspace{17\baselineskip}
According to what we have proved at the cited place, it therefore also follows that
\[
24.\quad\text{\textit{the integral}}\qquad
\int\frac{\mathrm{d}\varphi}{\mathord{\surd}{((Px)^4\cos\varphi^2+(Qx)^4\sin\varphi^2)}}
\]
\textit{extended from $\varphi=0$ to $\varphi=2\pi$, $=2\pi$, or also from $\varphi=0$ to $\varphi=\tfrac12k\pi$, where $k$ denotes any integer, $=\tfrac12k\pi$.}

\clearpage
\endgroup
